\section{Split bound} \label{sec:splitbound}

We now derive an upper bound on the weighted information gain
$\Gamma_T^w$ by exploiting spatial structure in the weighting function.
The regional transductive formulation discussed in
Appendix~\ref{app:transobjective} provides intuition for why decomposing
the domain into regions can be useful, while the independent-regions case
in Appendix~\ref{app:indepregions} shows that such decompositions become
exact when the Gaussian process does not couple different regions.

In the general setting, regions are correlated through the kernel and an
exact decomposition is no longer possible. Nevertheless, a similar idea
can be used to obtain an upper bound. The key observation is that posterior
variance can only decrease when conditioning on additional observations.
Consequently, the contribution from a region can be bounded by considering
only the subsequence of observations lying inside that region.

This leads to a decomposition of $\Gamma_T^w$ into regional information
gain terms with rescaled effective noise variances. Intuitively, larger
weights correspond to smaller effective noise levels, and therefore to
larger regional information gains.

\subsection{Statement of the split bound}

\begin{proposition}[Split bound for piecewise constant weights] \label{prop: splitbound}
Let $\mathcal{X}=\bigcup_{i=1}^m A_i$ be a disjoint partition, and let
\begin{equation*}
w(x)
=
\sum_{i=1}^m
c_i\mathbf{1}_{A_i}(x).
\end{equation*}
Then the weighted information gain satisfies
\begin{equation*}
\Gamma_T^w
\leq
\max_{T_1+\cdots+T_m=T}
\sum_{i=1}^m
\Gamma_{A_i}^{\mathrm{unweighted}}
\left(
T_i;\frac{\sigma^2}{c_i^2}
\right),
\end{equation*}
where $\Gamma_{A_i}^{\mathrm{unweighted}}(T_i;\sigma^2/c_i^2)$ denotes the ordinary information gain restricted to $A_i$ with effective noise variance $\sigma^2/c_i^2$.
\end{proposition}

\begin{remark}[Effective noise interpretation]
The split bound may be interpreted as decomposing the weighted information gain into regional information gain problems with rescaled noise levels. Larger weights correspond to smaller effective noise variances and therefore larger regional information gains.
\end{remark}

\subsection{Proof in the two-region case}

\begin{proof}[Proof of Proposition \ref{prop: splitbound}, two-region case]
We first prove the claim in the simple case where the ambient space is partitioned into two regions. Let
\begin{equation*}
w(z)
=
a\,\mathbf{1}_A(z)
+
b\,\mathbf{1}_B(z),
\end{equation*}
where
\begin{equation*}
A\cap B=\varnothing,
\qquad
A\cup B=\mathcal{X}.
\end{equation*}
Define
\begin{equation*}
\Gamma_A(N;a)
:= \frac{1}{2}
\max_{x_1,\dots,x_N\in A}
\sum_{j=1}^N
\log\!\left(
1+
\frac{a^2}{\sigma^2}
\tilde{\sigma}_{j-1,A}^2(x_j)
\right),
\end{equation*}
where $\tilde{\sigma}_{j-1,A}^2(x)$ denotes the posterior variance obtained by conditioning only on the previously selected points lying in $A$ (rather than on the full observation sequence)
and similarly
\begin{equation*}
\Gamma_B(N;b)
:= \frac{1}{2}
\max_{x_1,\dots,x_N\in B}
\sum_{j=1}^N
\log\!\left(
1+
\frac{b^2}{\sigma^2}
\tilde{\sigma}_{j-1,B}^2(x_j)
\right).
\end{equation*}
Take a sequence $z_1,\dots,z_N,$
and suppose there are $m$ points in $A$ and $N-m$ points in $B$.
Write the $A$-subsequence as $z_{t_1},\dots,z_{t_m}$
in order of appearance in the full sequence. Then its contribution is
\begin{equation*}
\sum_{j=1}^m
\log\!\left(
1+
\frac{a^2}{\sigma^2}
\tilde{\sigma}_{t_j-1,\mathcal{X}}^2(z_{t_j})
\right).
\end{equation*}
Now since $\tilde{\sigma}_{t_j-1,\mathcal{X}}^2(z_{t_j})
\leq
\tilde{\sigma}_{j-1,A}^2(z_{t_j}),$
(as conditioning on additional points can only reduce posterior variance),
\begin{equation*} \frac{1}{2}
\sum_{j=1}^m
\log\!\left(
1+
\frac{a^2}{\sigma^2}
\tilde{\sigma}_{t_j-1,\mathcal{X}}^2(z_{t_j})
\right)
\leq \frac{1}{2}
\sum_{j=1}^m
\log\!\left(
1+
\frac{a^2}{\sigma^2}
\tilde{\sigma}_{j-1,A}^2(z_{t_j})
\right)
\leq
\Gamma_A(m;a).
\end{equation*}
The same argument applies to the $B$-subsequence, giving
\begin{equation*}
\Gamma_N(w,k)
\leq
\Gamma_A(m;a)
+
\Gamma_B(N-m;b).
\end{equation*}
Taking the maximum over all allocations $m$ yields
\begin{equation*}
\Gamma_N(w,k)
\leq
\max_{m=0,\dots,N}
\left\{
\Gamma_A(m;a)
+
\Gamma_B(N-m;b)
\right\}.
\end{equation*}
Now observe that
\begin{equation*}
\Gamma_A(m;a)
=
\Gamma_A^{\mathrm{unweighted}}
\!\left(
m;
\left(\frac{\sigma}{a}\right)^2
\right),
\end{equation*}
corresponding to ordinary information gain on $A$ with effective noise variance
$\sigma_A^2
=
\frac{\sigma^2}{a^2}.$
Similarly,
\begin{equation*}
\Gamma_B(N-m;b)
=
\Gamma_B^{\mathrm{unweighted}}
\!\left(
N-m;
\left(\frac{\sigma}{b}\right)^2
\right).
\end{equation*}
Hence
\begin{equation*}
\Gamma_N(w,k)
\leq
\max_{m=0,\dots,N}
\left\{
\Gamma_A^{\mathrm{unweighted}}
\!\left(
m;
\left(\frac{\sigma}{a}\right)^2
\right)
+
\Gamma_B^{\mathrm{unweighted}}
\!\left(
N-m;
\left(\frac{\sigma}{b}\right)^2
\right)
\right\}.
\end{equation*}
Equivalently, in log-determinant form,
\begin{equation*}
\Gamma_A^{\mathrm{unweighted}}
\!\left(
m;
\left(\frac{\sigma}{a}\right)^2
\right)
= \frac{1}{2}
\max_{S\subseteq A,\ |S|=m}
\log\det
\left(
I+\sigma_A^{-2}K_{S,S}
\right),
\end{equation*}
so that
\begin{equation*}
\Gamma_N(w,k)
\leq \frac{1}{2}
\max_{m=0,\dots,N}
\Bigg\{
\max_{S\subseteq A,\ |S|=m}
\log\det
\left(
I+\left(\frac{\sigma}{a}\right)^{-2}
K_{S,S}
\right)
\end{equation*}

\begin{equation*}
\qquad\qquad
+
\max_{Y\subseteq B,\ |Y|=N-m}
\log\det
\left(
I+\left(\frac{\sigma}{b}\right)^{-2}
K_{Y,Y}
\right)
\Bigg\}.
\end{equation*}
\end{proof}

\subsection{Extension to finitely many regions}
\begin{proof}[Proof of Proposition \ref{prop: splitbound}, general case]
The same argument extends to finitely many disjoint regions. Let
\begin{equation*}
\mathcal{X}
=
\bigcup_{i=1}^m A_i,
\qquad
A_i\cap A_j=\varnothing
\quad
(i\neq j),
\end{equation*}
and define
\begin{equation*}
w(x)
=
\sum_{i=1}^m
c_i\mathbf{1}_{A_i}(x).
\end{equation*}
If $T_i$ denotes the number of sampled points in region $A_i$, with $T_1+\cdots+T_m=T,$
then applying the subsequence argument region-by-region gives
\begin{equation*}
\Gamma_T^w
\leq
\max_{T_1+\cdots+T_m=T}
\sum_{i=1}^m
\Gamma_{A_i}^{\mathrm{unweighted}}
\left(
T_i;\frac{\sigma^2}{c_i^2}
\right).
\end{equation*}
Equivalently, each region behaves like an independent information gain problem with effective noise level
$
\sigma_i^2
=
\frac{\sigma^2}{c_i^2}.$
\end{proof}

Unlike the exact decomposition in the independent-region setting, this bound remains valid even when correlations between regions are present. In the independent-region setting of Appendix \ref{app:indepregions}, the inequalities above become equalities, since observations outside a region no longer affect the posterior variances within it.


\subsection{Comparison with the crude bound}

The crude bound corresponds to replacing the weight function by its maximum value everywhere. If $a\geq b$, this gives
\begin{equation*}
B_{\mathrm{crude}}
:=
\Gamma_N^{\mathrm{unweighted}}
\left(k;\frac{\sigma^2}{a^2}\right)
= \frac{1}{2}
\max_{|S|=N}
\log\det
\left(
I+\frac{a^2}{\sigma^2}K_{S,S}
\right).
\end{equation*}
By contrast, the split bound from Proposition \ref{prop: splitbound} keeps track of how many samples fall in each region:
\begin{equation*}
B_{\mathrm{split}}
:=
\max_{m=0,\dots,N}
\left\{
\Gamma_A^{\mathrm{unweighted}}
\left(m;\frac{\sigma^2}{a^2}\right)
+
\Gamma_B^{\mathrm{unweighted}}
\left(N-m;\frac{\sigma^2}{b^2}\right)
\right\}.
\end{equation*}
The crude bound treats the whole domain as if it had the largest weight $a$, while the split bound only applies this effective noise level inside $A$ and uses the smaller effective weight in $B$.

With this decomposition, we hope to exploit geometric differences between regions. In particular, if highly weighted regions occupy smaller subsets of the domain, the associated information gain terms may decay more rapidly.

\subsection{Alternative weighting approaches}

The split bound exploits spatial structure in $w$ while leaving the
Gaussian process model unchanged. We contrast this with two alternative
ways of incorporating task relevance, discussed in detail in
Appendices~\ref{app:weighted_kernel} and
\ref{app:regional_objective}.

First, one could absorb the weights directly into the kernel by defining
$$
k_w(x,x')
=
w(x)w(x')k(x,x').
$$
For this modified kernel, the standard information gain machinery applies
directly. However, as discussed in Appendix~\ref{app:weighted_kernel}, this
changes the geometry of the Gaussian process prior rather than only changing
the acquisition objective. In particular, the posterior variances under the
weighted kernel do not generally satisfy
$$
\sigma_{t,w}^2(x)
=
w(x)^2\sigma_t^2(x).
$$
Therefore, kernel weighting corresponds to solving a different regression
problem: the notion of similarity between points is changed by the relevance
weights, rather than only changing which uncertainties are important.

Second, one can define a regional objective by applying weights to
transductive information gains over different target regions. Thus, the weighting is applied to the target region where uncertainty should be
reduced, rather than to the locations chosen for sampling. Given a
partition $\{A_i\}_{i=1}^m$ of set $S$], this gives
\[
F_w(S)
=
\sum_{i=1}^m c_i I(f_{A_i};y_S).
\]
Each term is submodular  and since nonnegative
sums of submodular functions remain submodular, the objective preserves
the usual greedy guarantees (see Appendix~\ref{app:regional_objective}).
However, this represents a different notion of task relevance: the
weights act on the prediction targets rather than directly on the
sequential variance reductions. 

The weighted acquisition objective studied here instead weights the marginal
value of reducing uncertainty at each sampled location. This allows a
graded notion of relevance across $\mathcal X$ while leaving the Gaussian
process posterior unchanged.
